% Section I — Introduction
\section{Introduction}\label{sec:intro}

Thermoreflectance in the time domain (TDTR) and in the frequency domain (FDTR) is a standard probe of
heat transport in thin films and across interfaces~\cite{Cahill2004,Schmidt2009}. When the heated volume or the
modulation period approaches phonon mean free paths or lifetimes, the data depart from the predictions
of Fourier's law: apparent conductivities depend on the modulation frequency~\cite{KohCahill2007,Regner2013}
and on the size and anisotropy of the heating~\cite{WilsonCahill2014}, and such departures have been
inverted for mean-free-path or relaxation-time spectra~\cite{Minnich2012,Forghani2016,Chen2025}. A second line
of analysis replaces Fourier's law by a constitutive law with a heat-flux relaxation
time~$\tau_R$---the Cattaneo law~\cite{JosephPreziosi1989}---or, in addition, a nonlocal
length~$\ell$---the Guyer--Krumhansl (GK) equation~\cite{GuyerKrumhansl1966a,GuyerKrumhansl1966b,GuoWang2015}---and
interprets thermoreflectance signals through these
parameters~\cite{Beardo2020,HennessyMyers2021,CamachoDeLaRosa2025}. Their values are then read
kinetically: $\tau_R$ as a resistive (momentum-relaxing) time and $\ell$ as a length set by normal
processes.

Aluminum nitride is a natural test case for this inference. Its thermal conductivity is high---about
$321~\mathrm{W\,m^{-1}K^{-1}}$ at room temperature was measured by TDTR on films grown to $18$ and
$22.5~\mu\mathrm{m}$~\cite{Cheng2020}, a value that cannot be transferred to a submicrometer film without
a measurement---and first-principles transport workflows for it have been benchmarked against one
another~\cite{McGaughey2025}; Callaway-type descriptions of its phonon transport have been examined
against first-principles rates~\cite{MaLiLuo2014}. A first-principles route from AlN to a
thin-film thermoreflectance prediction has three steps: (i)~static transport data---mode lifetimes
and the bulk conductivity tensor---are computed from harmonic and anharmonic force constants;
(ii)~an energy-conserving closure of the linearized Boltzmann equation is constructed from them;
(iii)~the closure enters a model of the film response. The closure coefficients that govern the
dynamics, above all the Cattaneo time, are not static quantities. This paper addresses the question:
\emph{what do static transport data determine about a conserving dynamic closure, and what further
operator or measurement information is needed?}

Several ingredients are established and are used here as attributed background. For the grey Callaway
model~\cite{Callaway1959}, the conserving closure has grad-div coefficient $1/3$ instead of the
coefficient $2$ of the original GK equation; Guo and Wang give it, including the combined time
$\tau_c$, in Eq.~(7) of Ref.~\cite{GuoWang2022}, Sendra \emph{et~al.} derive and discuss it~\cite{Sendra2021,Sendra2022},
and the coefficient $2$ has been attributed to an approximate inverse of the normal-process collision
matrix~\cite{HardyAlbers1974,GuoWang2022}; see also Ref.~\cite{Struchtrup2025}. The eigenvectors of the symmetrized collision operator
(relaxons) express the conductivity as a sum over relaxation modes~\cite{CepellottiMarzari2016}, and
conductivity and thermal viscosity are carried by relaxons of opposite parity~\cite{Simoncelli2020}.
Reductions of collision dynamics to a few collective variables are exact only under lumpability
conditions~\cite{WeiKuo1969,Aoki1968,LiRabitz1989,Buchholz1994} and otherwise generate memory terms of
projection type~\cite{Zwanzig1961,Mori1965}, on which a recent statistical theory of heat conduction with
spatiotemporal memory kernels is built~\cite{ZengDong2026}; descriptions built on diagonal rates alone can fail for
hydrodynamic transport~\cite{MalviyaRavichandran2023}. For Fourier films, Krapez and Rigollet showed that
thickness, diffusivity and conductivity of a coating are not separately identifiable from front-face
data, whatever the frequencies and noise level~\cite{KrapezRigollet2017}.

Within this background the paper makes the following contributions. Each is stated with its
hypotheses and one of the status labels defined below.
\begin{enumerate}
\item In entropy coordinates the Cattaneo time of a conserving closure, defined by matching a single-pole
law to the dynamic conductivity at zero frequency, is the first memory moment
$\tmem=-K'(0)/K(0)=|\Cop^+\jj|^2/\jj\tr\Cop^+\jj$ of the current response (Sec.~\ref{sec:kinetics});
in the grey Callaway model it equals $\tau_R$ exactly. This is a definition and an elementary check, not a new result.
\item Static data (lifetimes, invariants and the dc response) bound $\tmem$ from below by
$\tstat=K(0)/|\jj|^2$ (Proposition~\ref{prop:lower}) but, in the class of positive-semidefinite
operators, not from above (Theorem~\ref{thm:psd}). The mathematics is elementary: it is an instance of the
indeterminacy of a Stieltjes function from finitely many of its moments~\cite{Akhiezer2020,Schmudgen2017}.
We have not found an explicit identifiability statement of this kind for phonon kinetics in the
literature searched.
\item If the events are fixed by the harmonic spectrum and only their nonnegative rates are unknown,
$\tmem\le M^2K(0)$ with a sharp, rate-independent constant~$M$ (Theorem~\ref{thm:cone} and
Proposition~\ref{prop:sharp}). The inequality is a known consequence of the convex-hull theorem of Ben-Tal and
Teboulle for weighted least squares~\cite{BenTalTeboulle1990}; its application to the memory of kinetic
operators is potentially novel within the literature searched. A per-event lower bound on the rates bounds
$\tmem$ without any data (Proposition~\ref{prop:prior}).
\item In the relaxation-time approximation (RTA) the memory moment diverges in the continuum when the
infrared exponent of the acoustic rates satisfies $2\alpha\ge d$, while the conductivity is finite for
$\alpha<d$; for event operators the divergence follows under explicit hypotheses on the infrared
rates and solution (Theorem~\ref{thm:ir}), and then $M^2K(0)$ diverges under mesh refinement. No novelty is claimed
for this criterion; an analogous coexistence of a finite conductivity with divergent higher moments, of
the free-path distribution, underlies the superdiffusive description of quasiballistic transport in
alloys~\cite{Vermeersch2015a,Vermeersch2015b}.
\item For a one-dimensional front-face measurement on a Cattaneo or GK film on a substrate, the
depth-scaling gauge of the Fourier problem persists: the film enters only through
$(\xi_1,b,\tau_R,\tau_\ell)$, so neither $\ell$ nor the closure coefficient is separately identifiable
(Proposition~\ref{prop:gauge} and Corollary~\ref{cor:beta}). The Fourier gauge is
known~\cite{KrapezRigollet2017}; its extension to non-Fourier laws is elementary dimensional analysis and is
potentially novel within the literature searched. Whether the four combinations are themselves identifiable is
the subject of a theorem stated in Sec.~\ref{sec:frontface}: \PENDING{Theorem: generic rank four and degenerate set}.
\item \PENDING{AlN first-principles results: film-regularized memory-time ranges (Sec.~\ref{sec:aln})}
\PENDING{FDTR detectability of those ranges (Sec.~\ref{sec:fdtr})}
\end{enumerate}

The paper does not claim a value of the bulk memory time of AlN, an experimental validation, or any
part of the conserving closure itself. The Fourier resonance $\tau_R=\tau_\ell$ used in
Sec.~\ref{sec:frontface} was named by Both \emph{et~al.}~\cite{Both2016} and analyzed by
Kov\'acs~\cite{Kovacs2018}.

\emph{Status labels.} \status{proved}: a complete proof is given. \status{derived under assumptions}:
the derivation is complete given hypotheses that are stated but not proved (for instance a property of a
specific operator, or a model of the experiment). \status{open}: neither proved nor refuted. Numerical
results in Secs.~\ref{sec:aln} and \ref{sec:fdtr} will carry the label
\status{numerically demonstrated} together with their error estimates. Statements of novelty are
restricted to the literature we searched in October 2026: Crossref and OpenAlex (about fifty queries
each), arXiv full text, publisher and repository pages, and the forward citations of
Refs.~\cite{BenTalTeboulle1990,Stewart1989}.

Section~\ref{sec:kinetics} sets up the kinetic framework, the memory moment and the two operator
classes. Section~\ref{sec:identifiability} states and proves what static data do and do not determine.
Section~\ref{sec:frontface} formulates the front-face measurement model and its gauge.
Sections~\ref{sec:aln}--\ref{sec:conclusions} are reserved for the AlN results, the detectability
analysis and the conclusions. The notation is collected in Appendix~\ref{app:notation}.
